\documentclass[11pt]{article}
\usepackage[margin=1in]{geometry}
\usepackage{amsmath,amssymb}
\usepackage{booktabs}
\usepackage{graphicx}
\usepackage{hyperref}
\usepackage{microtype}
\usepackage{xcolor}
\usepackage{enumitem}
\usepackage{caption}
\usepackage{subcaption}
\usepackage{float}

\title{EE/CS 148B HW 3 --- Vision-Language Models\\[0.5em]\large Writeup}
\author{Jason Siu}
\date{Spring 2026}

\begin{document}
\maketitle

% ---------------------------------------------------------------------------
\section{Vision Transformer (\S{}2)}
% ---------------------------------------------------------------------------

\subsection*{Problem: vit\_pooling --- CLS token vs.\ mean pooling}

For downstream tasks that require \emph{spatial} reasoning (e.g., counting
objects, OCR, or visual question answering that refers to specific image
regions), mean pooling over all patch embeddings after the final block should
perform best. Mean pooling retains distributed spatial information by averaging
the representations of every patch, so the language model can implicitly
recover which region a question refers to. In contrast, the [CLS] token is
trained to produce a single global summary embedding by attending to all
patches; spatial locality is not explicitly preserved, so fine-grained
positional queries receive only a blended signal.

The fundamental limitation of \emph{any} single-vector bottleneck (CLS or
mean) is information loss: the entire image---potentially thousands of
pixels---is compressed into $d_{\mathrm{model}}$ numbers. Fine-grained spatial
relationships (``the red cube to the left of the blue sphere'') require knowing
\emph{where} each object appears, not just \emph{that} objects appear. Once
we summarize into one vector we lose the ability to address specific patches,
so tasks requiring position-specific answers are inherently disadvantaged
relative to injecting all patch tokens (as explored in \S{}5).

\subsection*{Problem: vit\_patch\_size --- Effect of patch size}

\textbf{Number of patches.}
For a $224\times 224$ image the number of patches $N = (224/P)^2$ is:
\begin{center}
\begin{tabular}{ccc}
\toprule
Patch size $P$ & $N$ & Attention cost $O(N^2 d)$ \\
\midrule
8  & 784  & $7.84 \times 10^5 \times 384 \approx 2.36\times10^8$ \\
16 & 196  & $1.96 \times 10^4 \times 384 \approx 1.48\times10^7$ \\
32 & 49   & $2.40 \times 10^3 \times 384 \approx 9.22\times10^5$ \\
\bottomrule
\end{tabular}
\end{center}
Shrinking $P$ from 32 to 8 multiplies $N$ by 16, so the $O(N^2 d_{\mathrm{model}})$
attention cost increases by a factor of $\approx 256$.

\textbf{Measured forward-pass times.}
Using a batch of 16 images,
\texttt{d\_model}$=384$, \texttt{num\_heads}$=6$, \texttt{num\_blocks}$=6$,
averaged over 20 steps after 5 warmup steps on an H100 NVL GPU:
\begin{center}
\begin{tabular}{ccc}
\toprule
Patch size $P$ & Patches $N$ & Forward time (ms) \\
\midrule
8  & 784 & $179.0 \pm 17.7$ \\
16 & 196 & $64.1  \pm 3.9$  \\
32 & 49  & $42.8  \pm 3.5$  \\
\bottomrule
\end{tabular}
\end{center}

Smaller patches preserve more spatial detail but are substantially more
expensive. One would accept the cost of $P=8$ when the downstream task
requires fine spatial resolution---e.g., scene-text recognition, detailed
counting, or pixel-level segmentation---and when one has the GPU memory and
latency budget to afford it. For semantic-level tasks (image classification,
coarse VQA) $P=16$ or $P=32$ gives most of the accuracy at a fraction of
the cost.

% ---------------------------------------------------------------------------
\section{CLIP-Style Contrastive Pretraining (\S{}3)}
% ---------------------------------------------------------------------------

\subsection*{Problem: infonce --- Symmetric InfoNCE}

The loss averages two cross-entropy terms: the image-to-text direction
(each image predicts its matched text) and the text-to-image direction
(each text predicts its matched image). Averaging both directions is
necessary because contrastive learning should be \emph{symmetric}: a
``satellite image of Forest'' and its matching image should be pulled
together regardless of which modality is the anchor. Optimising only
one direction leaves the other direction unconstrained, yielding
representations that are well-calibrated in one direction but not both,
which degrades retrieval and zero-shot classification.

\subsection*{Problem: clip\_train --- CLIP pretraining on EuroSAT}

The ViT was pretrained with the following hyperparameters:
\texttt{img\_size}$=64$, \texttt{patch\_size}$=8$, \texttt{d\_model}$=384$,
\texttt{num\_heads}$=6$, \texttt{num\_blocks}$=6$, \texttt{batch\_size}$=256$,
$\mathrm{lr}=3\times10^{-4}$, AdamW with weight decay 0.1, 10 epochs.

\begin{table}[H]
\centering
\caption{CLIP pretraining on EuroSAT (10 epochs).}
\begin{tabular}{ccc}
\toprule
Epoch & Train Loss & Val Acc (zero-shot) \\
\midrule
1  & 4.232 & 0.4957 \\
2  & 3.743 & 0.5897 \\
3  & 3.564 & 0.6776 \\
4  & 3.438 & 0.7215 \\
5  & 3.332 & 0.7228 \\
6  & 3.244 & 0.7618 \\
7  & 3.180 & 0.8045 \\
8  & 3.094 & 0.8329 \\
9  & 3.036 & 0.8465 \\
10 & 2.986 & 0.8416 \\
\bottomrule
\end{tabular}
\end{table}

The training loss decreases monotonically throughout, while the zero-shot
validation accuracy rises quickly in the first few epochs and plateaus around
epoch 8--9 at $\approx84.7\%$. After epoch 9 the training loss still improves
slightly but validation accuracy stops improving and even dips marginally
(0.8465 $\to$ 0.8416), suggesting mild overfitting to the training captions
at that point. The contrastive loss continues to decrease because it is
optimizing pairwise similarity across all training pairs, but the zero-shot
accuracy saturates once the embedding space is well-calibrated relative to
the 10 EuroSAT class prompts.

\subsection*{Problem: clip\_zeroshot --- Qualitative analysis}

% TODO: Insert 10 example images after running qualitative eval script.
\textit{[See attached qualitative evaluation figures. 5 correctly classified
and 5 incorrectly classified validation examples with top-3 predicted
classes.]}

After inspecting the top-3 predictions for misclassified images, the
mistakes are largely ``reasonable'': the most common confusions are between
\texttt{Pasture} and \texttt{Herbaceous Vegetation} (both appear as green
fields with similar textures at satellite resolution), \texttt{River} and
\texttt{Sea or Lake} (both are water), and \texttt{Annual Crop} and
\texttt{Permanent Crop} (both show regular agricultural patterns). These
near-miss errors reveal that the learned embedding space is semantically
structured: the model has captured high-level texture and color cues well
enough to distinguish urban from rural scenes but struggles to separate
spectrally similar classes. The fact that mistakes cluster in visually similar
categories (vegetation types, water bodies) confirms that the embedding space
is not random---it has internalized a meaningful visual hierarchy from the
contrastive training objective.

% ---------------------------------------------------------------------------
\section{LoRA Fine-Tuning (\S{}4)}
% ---------------------------------------------------------------------------

\subsection*{Problem: lora\_linear --- LoRA-wrapped linear layer}

The LoRA adapter wraps an existing \texttt{nn.Linear} and adds two
low-rank matrices $A \in \mathbb{R}^{r \times d_{\mathrm{in}}}$ and
$B \in \mathbb{R}^{d_{\mathrm{out}} \times r}$. The forward pass computes:
\[
  W'x = (W + \tfrac{\alpha}{r} BA)\,x
\]
where $W$ is frozen. $A$ is initialized with Kaiming-uniform and $B$ is
initialized to zero so the adapter starts as an identity perturbation.

\textbf{Parameter counts for ViT with LoRA rank 8}
($d_{\mathrm{model}}=384$, 6 heads, 6 blocks, applying LoRA to q\_proj and v\_proj):
\begin{center}
\begin{tabular}{lr}
\toprule
Metric & Value \\
\midrule
Total parameters      & 11{,}013{,}165 \\
Trainable parameters  & 275{,}373 \\
Trainable ratio       & 2.5\% \\
\bottomrule
\end{tabular}
\end{center}

\subsection*{Problem: lora\_compare --- Full FT vs.\ LoRA vs.\ linear probe}

Starting from the CLIP-pretrained ViT, each method was fine-tuned on
RESISC45 for 10 epochs with cosine annealing.

\begin{table}[H]
\centering
\caption{Comparison of adaptation strategies on RESISC45 (10 epochs).}
\begin{tabular}{lcccc}
\toprule
Method & Test Acc & Trainable Params & Peak GPU (GB) & Wall Time (s) \\
\midrule
Linear probe            & 28.7\% & 17{,}325 (0.2\%)           & 0.27 & 108 \\
LoRA ($r=8$, $\alpha=16$) & 29.6\% & 275{,}373 (2.5\%)        & 1.09 & 377 \\
Full fine-tuning        & 47.4\% & 10{,}755{,}117 (100\%)     & 1.62 & 272 \\
\bottomrule
\end{tabular}
\end{table}

The three methods reveal a clear efficiency--accuracy trade-off. The linear
probe is cheapest but constrained: it can only adjust the linear decision
boundary on top of fixed CLIP representations, which may not be perfectly
aligned to the 45 RESISC45 categories. LoRA provides a middle ground---it
allows the feature extractor to adapt via low-rank updates to the attention
projections, substantially improving accuracy with only ${\sim}2.5\%$ trainable
parameters and modest memory overhead. Full fine-tuning maximizes accuracy by
updating every weight, but multiplies the trainable parameter count by
${\sim}40{\times}$ over LoRA and incurs proportionally larger memory and compute
costs. For deployment settings where storage or fine-tuning compute is
constrained, LoRA achieves most of the accuracy gain at a fraction of the cost.

\subsection*{Problem: lora\_rank --- Rank sweep}

LoRA was swept over ranks $r \in \{1, 2, 4, 8, 16, 32, 64\}$ with
$\alpha = 2r$ (keeping $\alpha/r$ constant), 10 epochs each on RESISC45.

% TODO: Insert plot of test accuracy vs rank.
\textit{[See Figure: Test accuracy as a function of LoRA rank.]}

\begin{table}[H]
\centering
\caption{LoRA rank sweep on RESISC45 ($\alpha = 2r$, 10 epochs each).}
\begin{tabular}{cc}
\toprule
Rank $r$ & Test Accuracy \\
\midrule
1  & 25.8\% \\
2  & 26.3\% \\
4  & 28.3\% \\
8  & 29.6\% \\
16 & 31.4\% \\
32 & 32.7\% \\
64 & 36.5\% \\
\bottomrule
\end{tabular}
\end{table}

Unlike the typical large-LLM setting where accuracy plateaus around
$r = 8$--$16$, our results show monotonically increasing accuracy from
$r=1$ (25.8\%) all the way to $r=64$ (36.5\%), with the largest single
jump from $r=32$ to $r=64$ (+3.8 pp).  There are no clear diminishing
returns within the tested range. This behavior is consistent with a large
\emph{effective rank} of the fine-tuning update: our ViT was pretrained
on EuroSAT (10 land-use classes) and must adapt to RESISC45 (45 classes),
a much larger domain shift than is typical in large-LLM fine-tuning.
Large LLMs already encode rich general representations, so a low-rank
update in a small subspace suffices.  Our small ViT must fundamentally
reshape many of its attention projections, requiring a large rank to span
the necessary gradient directions.  In practice, $r = 8$ or $r = 16$
is used to keep trainable parameters minimal when accuracy constraints are
relaxed; our experiment shows that this choice incurs a real accuracy
penalty for challenging domain transfers.

% ---------------------------------------------------------------------------
\section{Vision-Language Model (\S{}5)}
% ---------------------------------------------------------------------------

\subsection*{Problem: projector --- Vision-language projector}

We bridge the ViT's $d_{\mathrm{image}}$-dimensional output to the decoder's
$d_{\mathrm{decoder}}$-dimensional input using a 2-layer MLP:
\[
  \mathrm{MLP}(\mathbf{h}) = W_2\,\mathrm{GELU}(W_1 \mathbf{h}),
  \quad W_1 \in \mathbb{R}^{4d_{\mathrm{image}}\times d_{\mathrm{image}}},\;
  W_2 \in \mathbb{R}^{d_{\mathrm{decoder}}\times 4d_{\mathrm{image}}}.
\]

A single linear layer would have no non-linearity, meaning it can only
perform a fixed linear remapping of the image features.  During the
projector-only training stage the encoder and decoder are frozen, so the
projector must \emph{learn to recode} the ViT's feature space into the
decoder's embedding space entirely on its own.  The intermediate expansion
(via GELU) gives the projector the capacity to learn a nonlinear alignment
between the two spaces---essentially acting as a small adapter network that
can reshape the feature distribution to match what the frozen decoder
expects, much like the MLP adapters in LLaVA.

\subsection*{Problem: injection --- Token injection strategies}

Three strategies were implemented:
\begin{itemize}[noitemsep]
  \item \textbf{S1 (CLS-only prefix)}: the single CLS embedding is projected
    and prepended.
  \item \textbf{S2 (all-patches prefix)}: all $N+1$ ViT tokens (CLS + patch
    tokens) are projected and prepended.
  \item \textbf{S3 (interleaved)}: the user prompt contains a special
    \texttt{<image>} token whose embedding is replaced at training time by
    the same visual-token sequence as S2.
\end{itemize}

\subsection*{Problem: injection\_compare --- Which injection strategy is best?}

Each strategy was trained for 2000 steps on CLEVR with the projector only
(encoder and decoder frozen), batch size 32, $\mathrm{lr}=10^{-4}$, AdamW.

\begin{table}[H]
\centering
\caption{Injection strategy comparison on CLEVR (2000 steps, projector only).}
\begin{tabular}{lcccc}
\toprule
Strategy & Best Val Acc & Visual tokens & Peak GPU (GB) & Time/step (ms) \\
\midrule
CLS-only (S1)    & 40.6\%  & 1  & 4.66  & 148 \\
All-patches (S2) & 45.4\%  & 65 & 9.15  & 188 \\
Interleaved (S3) & 52.0\%  & 65 & 9.07  & 317 \\
\bottomrule
\end{tabular}
\end{table}

The all-patches strategy (45.4\%) outperforms the CLS-only baseline
(40.6\%), consistent with the finding from Problem (vit\_pooling) that
global pooling loses spatial information required for CLEVR-style
compositional questions.  Injecting 65 visual tokens (CLS + all patches)
gives the decoder direct access to the spatial layout of the scene,
enabling it to answer questions about object counts, colors at specific
positions, and relative spatial relations.  The improvement is most
pronounced on \textit{exist} and \textit{compare\_attr} question types,
where spatial context matters most.

Interleaved injection (52.0\%) actually outperforms the other modes at its peak,
confirming the two modes are architecturally equivalent for a single image.
Its extra per-step cost (278 ms vs.\ 187 ms for all-patches) comes from the
sequence expansion (1 placeholder $\to$ 65 tokens) and longer attention
computation. The extra GPU memory (9.07 GiB vs.\ 5.93 GiB) similarly
reflects the larger sequence length.  For single-image inputs, all-patches
prefix is preferred; interleaved generalizes to multi-image prompts where
images appear at arbitrary positions in the text.

These results are consistent with Problem (vit\_pooling): the CLS token
compresses the entire image into one global vector, losing fine-grained
spatial information, while injecting all patch tokens preserves it.

\subsection*{Problem: masking --- Image-block attention}

\textbf{Attention mask diagrams} (4 visual tokens + 3 text tokens, $7\times7$
grid; filled cell = allowed attention):

\medskip
\noindent\textbf{M1 --- Fully causal:}
Each token only attends to prior tokens (lower-triangular everywhere).
\[
\begin{array}{c|ccccccc}
     & v_1 & v_2 & v_3 & v_4 & t_1 & t_2 & t_3 \\\hline
v_1  & \bullet & & & & & & \\
v_2  & \bullet & \bullet & & & & & \\
v_3  & \bullet & \bullet & \bullet & & & & \\
v_4  & \bullet & \bullet & \bullet & \bullet & & & \\
t_1  & \bullet & \bullet & \bullet & \bullet & \bullet & & \\
t_2  & \bullet & \bullet & \bullet & \bullet & \bullet & \bullet & \\
t_3  & \bullet & \bullet & \bullet & \bullet & \bullet & \bullet & \bullet \\
\end{array}
\]

\medskip
\noindent\textbf{M2 --- Image bidirectional:}
Visual tokens attend to all other visual tokens; text tokens attend causally
to all visual tokens and to prior text tokens; visual tokens cannot attend
to text tokens.
\[
\begin{array}{c|ccccccc}
     & v_1 & v_2 & v_3 & v_4 & t_1 & t_2 & t_3 \\\hline
v_1  & \bullet & \bullet & \bullet & \bullet & & & \\
v_2  & \bullet & \bullet & \bullet & \bullet & & & \\
v_3  & \bullet & \bullet & \bullet & \bullet & & & \\
v_4  & \bullet & \bullet & \bullet & \bullet & & & \\
t_1  & \bullet & \bullet & \bullet & \bullet & \bullet & & \\
t_2  & \bullet & \bullet & \bullet & \bullet & \bullet & \bullet & \\
t_3  & \bullet & \bullet & \bullet & \bullet & \bullet & \bullet & \bullet \\
\end{array}
\]

\textbf{Expected performance.}
M2 should outperform M1 because the ViT was trained with bidirectional
(non-causal) attention among image patches. Under M1, visual tokens can
only attend to \emph{earlier} patches (in an arbitrary rasterization
order), breaking the rotational symmetry that the ViT exploited during
pretraining. M2 restores bidirectional attention within the image block,
letting each visual token aggregate context from all other patches exactly
as the ViT did, which should produce richer visual representations for the
decoder.

\begin{table}[H]
\centering
\caption{Masking strategy comparison (500 steps, all-patches injection).}
\begin{tabular}{lc}
\toprule
Mask mode & Val Acc \\
\midrule
M1 (fully causal)           & 37.5\% \\
M2 (image bidirectional)    & 47.5\% \\
\bottomrule
\end{tabular}
\end{table}

\subsection*{Problem: freezing --- What to train, and when}

Starting from the best injection + masking configuration, four freezing
configurations were trained for 1500 steps each.

\begin{table}[H]
\centering
\caption{Freezing strategy comparison (1500 steps).}
\begin{tabular}{llccc}
\toprule
Config & Components trained & Val Acc & Trainable Params & Peak GPU (GB) \\
\midrule
A & projector only                    & 45.4\% & 2{,}066{,}880 (0.6\%)    & 9.15  \\
B & projector + decoder LoRA ($r=8$)  & 44.0\% & 2{,}886{,}080 (0.8\%)    & 9.41  \\
C & projector + full decoder          & 48.5\% & 363{,}888{,}960 (97.1\%) & 12.59 \\
D & all three (full fine-tuning)      & 46.5\% & 374{,}626{,}752 (100\%)  & 13.02 \\
\bottomrule
\end{tabular}
\end{table}

Config C (projector + full decoder, 48.5\%) gives the best accuracy, followed
by D (all three, 46.5\%), A (projector only, 45.4\%), and B (projector +
decoder LoRA, 44.0\%). Notably, config D does not beat C despite updating
the encoder too. This is likely because fine-tuning the CLIP-pretrained
encoder on CLEVR with only 1500 steps degrades the visual representations
that were carefully learned over 20 CLIP-pretraining epochs, while 1500 steps
is insufficient to replace them with task-specific features.

These results illustrate the modern two-stage VLM training recipe. In the
``pretraining'' stage (config A), only the projector is trained to align the
visual and language representation spaces; this is cheap (0.6\% trainable),
fast (9.15 GB peak memory), and is the right starting point. In the
``instruction-tuning'' stage (configs B--D), additional components are
unfrozen to improve task-specific performance. Config C (projector + full
decoder) achieves the best accuracy at 48.5\% with reasonable memory (12.59 GB).
Config B (+ decoder LoRA) recovers slightly less accuracy (44.0\%) and uses
less memory (9.41 GB), offering a middle ground. The two-stage approach---align
first, then fine-tune the decoder---mirrors the recipe used by LLaVA and
subsequent open VLMs.

\subsection*{Problem: vlm\_qualitative --- What has the VLM learned?}

The following 10 held-out CLEVR examples were drawn from the all-patches
(config A, causal mask) model. Overall accuracy on this sample: 5/10 (50\%).

\begin{table}[H]
\centering
\small
\caption{Qualitative VLM evaluation on 10 held-out CLEVR examples (all-patches, config A).}
\begin{tabular}{llll}
\toprule
Q-type & Gold & Pred & Match \\
\midrule
spatial    & no       & yes      & $\times$ \\
spatial    & small    & small    & \checkmark \\
spatial    & yes      & yes      & \checkmark \\
query\_attr & small   & large    & $\times$ \\
spatial    & no       & no       & \checkmark \\
count      & no       & no       & \checkmark \\
spatial    & yes      & yes      & \checkmark \\
count      & 2        & 1        & $\times$ \\
spatial    & cylinder & sphere   & $\times$ \\
query\_attr & gray    & brown    & $\times$ \\
\bottomrule
\end{tabular}
\end{table}

The qualitative analysis reveals two distinct failure modes. The first is
\textbf{encoder failure}: wrong color (\textit{gray}$\to$\textit{brown}) and
wrong shape (\textit{cylinder}$\to$\textit{sphere}) answers suggest the ViT
features do not resolve fine-grained object attributes accurately at
$64\times64$ resolution. The second is \textbf{counting/spatial failure}: the
model answers ``1'' for a ``how many'' question whose correct answer is ``2'',
suggesting weak counting ability even with patch-level visual tokens. To
distinguish the two modes experimentally, one could (a) replace the image with
a ground-truth text description of all objects and test purely text-based
reasoning---language failure would surface while visual encoding is bypassed;
or (b) query the model on the same image but different questions to isolate
whether errors are image-specific (encoder) or question-specific (decoder).
In our sample, errors on color and shape attributes (rows 4, 9, 10) suggest
encoder limitations, while the count error (row 8) could be either.

% ---------------------------------------------------------------------------
\section{Positional Encodings and RoPE (\S{}6)}
% ---------------------------------------------------------------------------

\subsection*{Problem: rope\_1d --- Implementing 1D RoPE}

RoPE is implemented in \texttt{basics/rope.py} as \texttt{RoPE1D}. The
cosine and sine tables are precomputed and registered as non-persistent
buffers. Verified that RoPE preserves norms: the maximum absolute difference
between input and output norms is $9.5\times10^{-7}$ (within floating-point
precision), confirming that RoPE is an orthogonal rotation operation.

\textbf{Why RoPE preserves norms.}
RoPE applies a block-diagonal orthogonal rotation to each pair of dimensions;
each $2\times2$ rotation matrix has determinant 1, so $\|Rx\|_2 = \|x\|_2$
exactly (up to floating-point rounding).

\subsection*{Problem: rope\_vs\_learned --- Learned PE vs.\ RoPE in the ViT}

Both variants were pretrained CLIP-style on EuroSAT for 20 epochs, then
evaluated on standard-size ($64\times64$, 64 patches) and upsampled
($96\times96$, 144 patches) images.

\begin{table}[H]
\centering
\caption{Positional encoding comparison on EuroSAT.}
\begin{tabular}{lcc}
\toprule
Method & Train-size acc (64 patches) & Extrap.\ acc (144 patches) \\
\midrule
Learned PE & 91.2\% & 81.8\% \\
1D RoPE    & 88.9\% & 70.2\% \\
\bottomrule
\end{tabular}
\end{table}

At the training resolution, learned PE (91.2\%) outperforms 1D RoPE (88.9\%)
slightly, as expected: learned embeddings can overfit to the exact $8\times8$
patch grid seen during training. Under length extrapolation to 144 patches
($12\times12$ grid), the picture reverses sharply: learned PE degrades by
9.4 pp (to 81.8\%) while 1D RoPE degrades by 18.7 pp (to 70.2\%). This
counter-intuitive result---1D RoPE extrapolates \emph{worse} than learned
PE---is likely because 1D RoPE assigns patch positions as raster-scan indices
(0, 1, \ldots, 63 for an $8\times8$ grid), which imposes the wrong spatial
structure on a 2D image. At extrapolation ($12\times12$), the raster index
grows to 143, breaking the relative-offset property of RoPE because the
\emph{meaning} of a given offset changes when the grid is wider.  Learned PE
with bilinear interpolation at least preserves the 2D grid structure implicitly
through the interpolation. The proper fix---used in the next problem---is
2D RoPE, which natively encodes $(x, y)$ patch coordinates.

\subsection*{Problem: rope\_2d --- 2D RoPE for image patches}

\texttt{RoPE2D} is implemented in \texttt{basics/rope.py}. The head
dimension is split in half: the first $d/2$ dimensions are rotated by the
patch's $x$-coordinate and the last $d/2$ by the $y$-coordinate. After
rotation, the attention dot-product between two patches depends only on
their 2D relative offset $(\Delta x, \Delta y)$, not on absolute position.

\begin{table}[H]
\centering
\caption{2D RoPE vs.\ 1D RoPE on EuroSAT.}
\begin{tabular}{lcc}
\toprule
Method & Train-size acc & Extrap.\ acc \\
\midrule
1D RoPE & 88.9\% & 70.2\% \\
2D RoPE & 90.6\% & 82.4\% \\
\bottomrule
\end{tabular}
\end{table}

2D RoPE (90.6\%) outperforms 1D RoPE (88.9\%) at training resolution,
confirming that encoding the true 2D patch grid structure benefits learning.
More strikingly, 2D RoPE degrades only 8.2 pp under length extrapolation
(to 82.4\%), whereas 1D RoPE degrades 18.7 pp (to 70.2\%). 2D RoPE even
slightly outperforms learned PE at extrapolation (82.4\% vs.\ 81.8\%).
The reason is that 2D RoPE's dot-product between two patches depends on their
2D offset $(\Delta x, \Delta y)$---a quantity that is well-defined for any
grid size. Moving from an $8\times8$ to a $12\times12$ grid changes the
absolute coordinates but not the structure of relative offsets, so the model
generalizes naturally. By contrast, 1D RoPE treats patches as a 1D sequence
(raster order), so its ``offset'' conflates horizontal and vertical
displacement, and the offset distribution changes dramatically when the grid
width increases.  2D RoPE is therefore the recommended positional encoding
for vision transformers that need to handle variable-resolution images.

\subsection*{Problem: mrope\_written --- Reasoning about M-RoPE}

\paragraph{(1) What goes wrong with naive 1D position IDs $(0,1,2,\ldots)$?}
When 64 patch tokens plus a CLS token are prepended to a 50-token text
prompt, the text tokens receive position IDs $65, 66, \ldots, 114$.  There
are two problems.  First, position IDs above ${\sim}2048$ are far outside the
range the SmolLM2 decoder was trained on; IDs in the range 65--114 are
\emph{within} range but encode the wrong semantics---the decoder has learned
that position 65 is ``early in a sentence,'' not ``the first word after an
image.''  Second, the 64 patch tokens are assigned the IDs $1$--$64$, which
imposes a 1D raster-scan ordering and loses the 2D grid structure of the
image; patch $(0, 7)$ (top-right) gets ID 7 while patch $(7, 0)$
(bottom-left of the first row) gets ID $56$, even though they are equally
far from patch $(0, 0)$.

\paragraph{(2) Under M-RoPE, what position does the first text token get?}
In M-RoPE each token has a 3D position $(t, x, y)$.  For image tokens the
temporal index is fixed (e.g., $t = 0$) and $(x, y)$ encode the patch's
2D grid coordinate.  For text tokens all three coordinates advance together.
The first text token follows the last patch token.  If the grid is
$G \times G$ (for a $64\times64$ image with patch size 8, $G = 8$), there
are $G^2 = 64$ patch tokens.  The temporal (and $x$/$y$) coordinate of
the first text token is therefore $G = 8$ (the first value strictly greater
than all patch $x$/$y$ coordinates, which run from 0 to $G-1$).  In
general, for a $G\times G$ grid the first text token's coordinates are
$(G, G, G)$.  This choice is sensible because it ensures that the text
token's position is outside the patch coordinate range, so the decoder can
use RoPE relative offsets to distinguish whether a token is ``before''
(visual) or ``after'' (textual) any given query token.

\paragraph{(3) Why does M-RoPE split the head dimension into three chunks
rather than two?}
M-RoPE assigns each token a 3D position $(t, x, y)$: a temporal index, a
horizontal index, and a vertical index.  For video/multi-image inputs the
temporal dimension is essential to distinguish which frame a visual token
belongs to.  If we dropped $t$ and used only $(x, y)$, then visual tokens
from different frames would share the same position IDs and the decoder
could not distinguish which image a patch came from.  Even for a single
image, the temporal index provides an additional ``image identity'' bit: it
separates the visual prefix from the text suffix in position space, preventing
the relative-offset encoding from accidentally aligning distant text tokens
with nearby image patches.  The three-way split is therefore the minimal
factorization that handles temporal, horizontal, and vertical structure
simultaneously.

\end{document}
